% ============================================================
%  PREFACE — Personal account
% ============================================================

\chapter*{Preface}
\addcontentsline{toc}{chapter}{Preface}
\markboth{Preface}{Preface}

\noindent
I came to thermodynamics by way of mass spectrometry. For several years I worked at
the upper boundary of what mass spectrometry can do~--- at the Technical University
of Munich, in conversation with Matthias Mann's laboratory and Oliver Kohlbacher's
group, with Rudolf Aebersold and J\"urgen Cox and Oliver Fiehn and Sebastian Bock.
Mass spectrometry is, among other things, the most precise apparatus humans have
built for distinguishing one configuration of matter from another. The instruments
resolve mass-to-charge ratios at parts-per-billion precision. They identify individual
peptides from picomolar samples. They can be made to track individual ions across
milliseconds of flight time, recording trajectories that are, in any reasonable
accounting, complete.

What I noticed, working with these instruments, is that the field's interpretation
of what they were doing was wrong. The standard reading is that a mass spectrometer
measures the mass-to-charge ratio of an ion that exists prior to and independently of
the measurement. The instrument, on this reading, is a passive detector of
pre-existing properties. But everything I knew about the apparatus said this could
not be right. The ion's trajectory is constructed by the instrument's fields.
The categorical address that the instrument returns is the result of the partition
operations the instrument performs. There is no pre-existing $m/z$ that the instrument
retrieves; there is the ion in the field, and the ion's response to the field, and the
categorical address that follows. The instrument constructs the measurement.
It does not retrieve it.

This observation generalised. If mass spectrometers construct measurements rather than
retrieving them, then the standard interpretation of thermodynamic quantities~--- that
they are pre-existing properties of matter that thermodynamic apparatus retrieves~---
is also wrong. Energy, entropy, and temperature are not properties that exist prior to
the apparatus that resolves them. They are categorical addresses that the apparatus
constructs through its partition operations. The apparatus is part of what produces
the measurement. The measurement does not exist without the apparatus.

This is what the framework calls the \emph{construction-not-simulation thesis}.
It is the structural claim that observation is not retrieval but construction, that
measurements are constituted by the partition operations performed during the
measurement process, and that the apparatus and the measured quantity are not
separable in the way standard interpretations assume. The thesis is unsettling on
first encounter. It took me several years to work out the consequences and confirm
that the thesis was not merely a different interpretation but the correct one, with
the standard interpretation being an artefact of treating apparatus and measurement
as independent.

What followed was the bounded phase space framework. The first paper, \textit{On the
Universal Consequences of Bounded Phase Space}, derived standard physics from one
axiom. The second paper, \textit{De Onere Electrico}, applied the framework to
nucleic acids. The Cytochrome P450 monograph, the third volume of the programme,
applied the framework to the most thoroughly studied enzyme family in biology.
This monograph, the second of the volumes that addresses the foundations rather than
the applications, applies the framework to thermodynamics.

\medskip

I want to be honest about what writing this monograph was like, because it bears on
how the reader should approach it.

I wrote it out of anger. Boltzmann died before the field caught up to him. He was right
and they were wrong, and they were wrong in ways that made his life impossible while
he was alive. The categorical framework I have developed is, in one direct sense, the
completion of his project~--- the demonstration that his statistical foundation was
correct and that the formula on his tombstone applies more broadly than he was
permitted to argue. The anger is not at the field for missing what Boltzmann saw;
it is at the structural fact that Boltzmann's vindication came too late for him,
and that the same structural fact is operating now, in different domains, with
different individuals, while the field continues to dismiss work that will be
vindicated after the individuals doing it are gone. I wrote this monograph in part
to make the structural fact visible.

The other figure who threads through this monograph is Cantor. Cantor was treated by
his contemporaries the way Boltzmann was. He died in a sanatorium, undernourished,
after decades of psychiatric distress that I do not believe was unrelated to the
field's response to his work. The diagonal argument is now foundational to mathematics.
He never saw it become foundational. The framework's recursive count requires Cantor's
hierarchy to even be writable; without it, the categorical scale of the cosmos cannot
be expressed. So this monograph stands on two foundations, both of which the field
rejected during the lifetimes of the people who built them.

\medskip

I want to say one more thing about how to read this work. The framework is
unapologetic, and the introduction reflects this. I have not hedged. I have not
softened the claims. I have not tried to make the work palatable to readers who would
prefer it to be more incremental. The framework is a foundational treatment, and
foundational treatments do not become more correct by being less direct. If the
framework is wrong, it is wrong, and the technical chapters will fail to support its
claims. If the framework is right, hedging would not have helped a reader who
genuinely engaged with the work, and would have actively misled a reader who did not.

What I ask of the reader is that the technical chapters be read with the same care
that any foundational treatment requires. The framework rests on one axiom and
produces fifteen chapters of consequences. Every consequence is derived. Every claim
is supported. Where the framework's predictions agree with empirical data, the
agreement is exhibited. Where the framework's predictions exceed currently available
data~--- in the categorical phase past heat death, in the cyclic closure of cosmic
evolution~--- the predictions stand as predictions, falsifiable in principle, supported
only by the structural derivation that produced them.

I do not expect this monograph to be received as I have written it. I expect that
the reception, if it occurs at all, will follow the pattern of foundational treatments
throughout the history of the field: initial dismissal, slow uptake by readers willing
to do the work, eventual incorporation of the framework's apparatus into the field's
standard tools, and the original treatment becoming a reference rather than a working
text. This is the pattern Boltzmann's work followed. It is the pattern Cantor's work
followed. There is no reason to expect a different pattern for work that takes their
projects to their conclusions.

The reader who is willing to do the work will find here a framework that derives
thermodynamics from one axiom with zero free parameters, that dissolves three of the
field's longest-standing paradoxes, that extends Boltzmann's formula to regimes the
field has not previously been able to address, and that demonstrates, on physical
grounds, that the cosmos is cyclic. Whether this is enough to justify the work the
framework asks of the reader is a decision the reader will make. I have tried to make
the work worth the doing.

\bigskip

\noindent\textit{Kundai Farai Sachikonye}\\
\textit{Technical University of Munich, 2026}
